\documentclass[10pt,landscape]{article}
\usepackage[landscape,margin=0.5in]{geometry}
\usepackage{multicol}
\usepackage{array}
\usepackage{booktabs}
\usepackage{xcolor}
\usepackage{colortbl}
\usepackage{tikz}
\usepackage{fancyhdr}

\definecolor{rom}{HTML}{224488}
\definecolor{zp}{HTML}{448822}
\definecolor{dim}{HTML}{888888}

\pagestyle{empty}
\setlength{\parindent}{0pt}
\setlength{\parskip}{3pt}
\newcommand{\hex}[1]{\texttt{\$#1}}
\newcommand{\sect}[1]{\vspace{4pt}{\large\bfseries\color{rom}#1}\vspace{2pt}\hrule\vspace{3pt}}
\newcommand{\code}[1]{\texttt{\small #1}}

\begin{document}
\begin{multicols}{3}

\begin{center}
{\LARGE\bfseries 6502life ROM Cheat Sheet}\\[2pt]
{\small Geometry tables at \hex{E000}--\hex{EE3F}}
\end{center}

% ---- Cell Index Grid ----
\sect{Neighborhood Cell Indices}

49 cells in spiral order. Cell $n$ at address $n \times 1024$.

\smallskip
{\small\centering
\begin{tabular}{r|ccccccc}
      & \multicolumn{7}{c}{$x$} \\
 $y$  & $-3$ & $-2$ & $-1$ & 0 & $+1$ & $+2$ & $+3$ \\ \midrule
 $+3$ & 48 & 44 & 36 & 25 & 29 & 37 & 45 \\
 $+2$ & 43 & 24 & 20 &  9 & 13 & 21 & 38 \\
 $+1$ & 35 & 19 &  8 &  \cellcolor{zp!20}\textbf{1} &  \cellcolor{zp!20}\textbf{5} & 14 & 30 \\
   0  & 28 & 12 &  \cellcolor{zp!20}\textbf{4} &  \cellcolor{rom!20}\textbf{0} &  \cellcolor{zp!20}\textbf{2} & 10 & 26 \\
 $-1$ & 34 & 18 &  7 &  \cellcolor{zp!20}\textbf{3} &  \cellcolor{zp!20}\textbf{6} & 15 & 31 \\
 $-2$ & 42 & 23 & 17 & 11 & 16 & 22 & 39 \\
 $-3$ & 47 & 41 & 33 & 27 & 32 & 40 & 46 \\
\end{tabular}
}

\smallskip
{\footnotesize
\textbf{0}=self,
\textbf{1}=N$(0,\!+\!1)$,
\textbf{2}=E$(+\!1,\!0)$,
\textbf{3}=S$(0,\!-\!1)$,
\textbf{4}=W$(-\!1,\!0)$\\
\textbf{5}=NE, \textbf{6}=SE, \textbf{7}=SW, \textbf{8}=NW,
\textbf{9--12}=N$^2$E$^2$S$^2$W$^2$
}

% ---- Cell address ----
\sect{Cell Base Address}

\code{base = cell\_index * \$0400}\\
High byte = \code{cell\_index * 4} \quad (low byte always \hex{00})

\smallskip
{\small
\begin{tabular}{lll}
Cell 0 (self) & \hex{0000} & Cell 5 (NE) \hex{1400} \\
Cell 1 (N)    & \hex{0400} & Cell 6 (SE) \hex{1800} \\
Cell 2 (E)    & \hex{0800} & Cell 7 (SW) \hex{1C00} \\
Cell 3 (S)    & \hex{0C00} & Cell 8 (NW) \hex{2000} \\
Cell 4 (W)    & \hex{1000} & Cell 9 (N$^2$) \hex{2400} \\
\end{tabular}
}

\smallskip
{\footnotesize\textbf{At runtime:}
\code{ASL; ASL} to get high byte from cell index.\\
Store in ZP pointer, use \code{LDA (\$ptr),Y} for indirect access.}

% ---- ROM tables ----
\sect{ROM Table Map}

Each row = 64 bytes. \code{row\_addr = \$E000 + row * 64}.\\
Index by cell index (0--48). Out-of-range returns \hex{FF}.

\smallskip
{\small
\begin{tabular}{@{}lll@{}}
\toprule
Address & Row & Operation \\ \midrule
\hex{E000}+$64j$ & 0--48 & Translate by cell $j$ \\
\hex{EC40} & 49 & Rotate 90\textdegree{} CW \\
\hex{EC80} & 50 & Rotate 180\textdegree \\
\hex{ECC0} & 51 & Rotate 270\textdegree{} CW \\
\hex{ED00} & 52 & Reflect X (flip $y$) \\
\hex{ED40} & 53 & Reflect Y (flip $x$) \\
\hex{ED80} & 54 & $x$-coord $+\,3$ \\
\hex{EDC0} & 55 & $y$-coord $+\,3$ \\
\hex{EE00} & 56 & $(x,y) \to$ cell index \\
\bottomrule
\end{tabular}
}

% ---- Translation ----
\sect{Translation (Vector Addition)}

\code{translate[j][i]} = cell at $\vec{v}_i + \vec{v}_j$

\smallskip
{\footnotesize
\textbf{Example:} cell 1's east neighbor = cell at $(0,+1)+(+1,0)$:}

\code{LDX \#2; LDA \$E040,X} $\to$ \code{A=5} (NE)

\smallskip
{\footnotesize
\textbf{Use:} find neighbor-of-neighbor, compose directions,\\
check reachability (result $\neq$ \hex{FF} means in range).}

% ---- Rotation ----
\sect{Rotation \& Reflection}

\begin{tabular}{@{}lll@{}}
\code{LDA \$EC40,X} & Rot 90 CW  & N$\to$E$\to$S$\to$W \\
\code{LDA \$EC80,X} & Rot 180    & N$\to$S, E$\to$W \\
\code{LDA \$ECC0,X} & Rot 270 CW & N$\to$W$\to$S$\to$E \\
\code{LDA \$ED00,X} & Reflect X  & N$\leftrightarrow$S \\
\code{LDA \$ED40,X} & Reflect Y  & E$\leftrightarrow$W \\
\end{tabular}

\smallskip
{\footnotesize\textbf{Iterate cardinals:}}

\code{LDX \#1 / @loop: ... LDA \$EC40,X / TAX / CPX \#1 / BNE @loop}

% ---- Coordinates ----
\sect{Coordinate Tables}

\code{LDX \#\textit{cell}; LDA \$ED80,X} $\to$ $x + 3$ \quad(range 0--6)\\
\code{LDX \#\textit{cell}; LDA \$EDC0,X} $\to$ $y + 3$ \quad(range 0--6)

\smallskip
Origin (cell 0) = coordinate $(3,3)$.

\smallskip
{\footnotesize\textbf{Reverse:} \hex{EE00}\code{+}$(x+3)$\code{+}$(y+3)*8$ $\to$ cell index.}

\code{LDA \$EE14} $\to$ 6 \quad{\footnotesize (offset $(+1,-1)$: $x$\code{=4}, $y$\code{=2}, addr \hex{EE00}+4+16)}

% ---- Oriented Registers ----
\sect{Oriented Registers (\$F0--\$F8)}

Top 6 bits = cell index, auto-rotated by memory mapper.\\
Bottom 2 bits = user data (preserved across rotations).

\code{value = (cell\_index << 2) | flags}

\smallskip
{\footnotesize
Write a direction to \hex{F0}; after interrupt + re-orientation,\\
\hex{F0} still points to the \emph{same relative neighbor}.\\
No compass needed for rotation-invariant programs.

\smallskip
\textbf{Note:} \hex{F9} (PCHI) is also auto-rotated but is\\
reserved for the saved program counter high byte.}

% ---- Compass ----
\sect{Compass (\$FA)}

Requires \code{hasCompass=true}.
\hex{FA}\code{ = orientation << 2}.

\smallskip
\begin{tabular}{@{}cccc@{}}
\hex{00} = 0\textdegree &
\hex{04} = 90\textdegree &
\hex{08} = 180\textdegree &
\hex{0C} = 270\textdegree \\
\end{tabular}

\smallskip
To access absolute direction $d$ (S=1,E=2,N=3,W=4):\\
\code{local = ((d-1-orient+4) \% 4) + 1}

\smallskip
{\footnotesize\textbf{Lookup table} (16 bytes, indexed by \hex{FA}\code{+}$(d-1)$):}

{\ttfamily\footnotesize
\begin{tabular}{l@{\,}l}
orient 0: & 1 2 3 4 \\
orient 1: & 4 1 2 3 \\
orient 2: & 3 4 1 2 \\
orient 3: & 2 3 4 1 \\
\end{tabular}
}

\code{LDX \$FA; LDA tbl+2,X} $\to$ local index for abs.\ North

% ---- Zero Page Conventions ----
\sect{Zero Page Layout (cell 0)}

{\small
\begin{tabular}{@{}ll@{}}
\hex{00}--\hex{9F}  & Code / scratch \\
\hex{A0}            & Key input buffer \\
\hex{A1}            & Cell unique ID \\
\hex{A5}            & Type tag \\
\hex{A6}            & State length \\
\hex{A8}--\hex{B7}  & State string (16 chars) \\
\hex{F0}--\hex{F8}  & Oriented registers \\
\hex{F9}            & Saved PC high (oriented) \\
\hex{FA}            & Compass / saved PC low \\
\hex{FB}            & Saved P register \\
\hex{FC}--\hex{FF}  & RNG (refreshed each interrupt) \\
\end{tabular}
}

% ---- Indirect pattern ----
\sect{Runtime Cell Access Pattern}

{\footnotesize
When the target cell index is computed at runtime (e.g.\ from\\
compass resolution), use indirect addressing:}

\begin{verbatim}
  ASL             ; cell_index * 2
  ASL             ; * 4 = addr high byte
  STA $02         ; store high byte
  LDA #$00
  STA $01         ; low byte = 0
  LDY #$A5        ; TYPE_TAG_OFFSET
  LDA ($01),Y     ; read type tag
\end{verbatim}

% ---- Quick Reference ----
\sect{BRK Operations}

{\small
\begin{tabular}{@{}lll@{}}
\code{BRK \$00}    & Reset  & PC $\to$ 0, yield \\
\code{BRK \$01--30}& Swap   & cell 0 $\leftrightarrow$ cell $b$ \\
\code{BRK \$31--60}& Copy   & cell 0 $\to$ cell $(b{-}48)$ \\
\end{tabular}
}

\smallskip
{\footnotesize
Swap is $O(1)$ ($\sim$49 cyc). Copy is $O(M)$ ($\sim$6000 cyc).\\
Fails silently if remaining cycles $<$ cost.}

\end{multicols}
\end{document}
